\section{Overview of the proof}\label{sec:overview}
\subsection{From factors to fractional tilings}\label{sec:sketch-factor}
In Section~\ref{factor402:section}, we prove that Theorem~\ref{compact390:main} and Theorem~\ref{compact390:finite} are equivalent, as follows.
Given a digraph $G$ as in Theorem~\ref{compact390:main}, we first set aside a small vertex set $A$ which can absorb any small $r$-divisible set of leftover vertices into an $H$-factor (Lemma~\ref{general:absorb}).
We then apply the regularity lemma to $G-A$. Its reduced digraph $R$ has minimum semidegree greater than $|R|/2$, so Theorem~\ref{compact390:finite} gives a perfect fractional homomorphic $H$-tiling of $R$, and rounding this tiling gives vertex-disjoint copies of $H$ covering all but a few vertices of $G-A$ (Lemma~\ref{framework390:embed}). Finally, $A$ absorbs the remaining vertices.
Thus, the main work is to prove Theorem~\ref{compact390:finite}, and for the rest of this overview we fix an orientation $H$ of an even cycle of length $r\geq 4$ and a digraph $F$ on $N$ vertices with $\delta^0(F)>N/2$.

\subsection{Notation}\label{sec:notation}
For a vertex $v$ of a digraph $G$, we write $N^+(v)$ and $N^-(v)$ for its outneighbourhood and inneighbourhood, respectively, and let $d^\pm(v)=|N^\pm(v)|$. Where we use $\pm$ or $\sigma\in\{+,-\}$ in an expression, we mean that it holds for both signs.
We write $xy$, or $x\to y$, for an edge directed from $x$ to $y$, and $x\leftrightarrow y$ if both $xy$ and $yx$ are edges, in which case we call $\{x,y\}$ a \emph{$2$-cycle}.
Unless stated otherwise, paths, walks and cycles in a digraph are directed, so that their edges are oriented consistently. A walk may repeat vertices, and a \emph{closed $\ell$-walk} is a closed walk with $\ell$ edges.
For a weight function $w$ on the vertices of a digraph and a vertex set $A$, we write $w(A)=\sum_{v\in A}w(v)$, and the \emph{cost} of a homomorphism or walk is the sum of the weights of its vertices, counted with repetition.
The digraph $D_{p,q}$ consists of two internally vertex-disjoint paths of lengths $p$ and $q$ with the same initial vertex and the same final vertex, and $Q$ is the orientation of a four-cycle with edges $ab,bc,cd,ad$, that is, $Q=D_{3,1}$.
The \emph{converse} $H^{\rm op}$ of a digraph $H$ is obtained by reversing all its edges. Replacing $H$ and $F$ by their converses preserves minimum semidegrees, homomorphisms and tilings, and we use this symmetry repeatedly.

We describe an orientation $H$ of a cycle by reading its vertices in a cyclic order and recording, for each edge, the sign $+$ if the edge is directed forwards and $-$ if it is directed backwards. This gives a cyclic word in $+$ and $-$; for example, $Q$ corresponds to $+++-$. We refer to the vertices of $H$ as the \emph{positions} of this word. The partial sums of the word form a walk on the integers with steps $\pm1$, which we call the \emph{height walk} of the word, and the \emph{height} of a position is the corresponding partial sum. The \emph{span} of the walk is its maximum height minus its minimum height, and it is \emph{closed} if it ends at its initial height.
A position is a \emph{turn} if its two incident edges point in opposite directions along the reading (so that both point into it or both point out of it), and it is \emph{internal} otherwise. A maximal set of consecutive edges with the same sign is a \emph{directed run}.

\subsection{Covers, and why the semidegree condition is natural}\label{sec:sketch-covers}
A \emph{fractional $H$-cover} of $F$ is a function $w:V(F)\to\mathbb R_{\ge0}$ such that every homomorphism from $H$ to $F$ has cost at least one.
By linear programming duality (Lemma~\ref{general:LP}), $F$ has a perfect fractional homomorphic $H$-tiling if and only if every fractional $H$-cover of $F$ has total weight at least $N/r$. So, suppose that $w$ is an optimal fractional $H$-cover with total weight $W<N/r$. We will derive a contradiction.

To see how the condition $\delta^0(F)>N/2$ can be used, delete one position $j$ from $H$, and consider a homomorphism $Q$ from the resulting path into $F$, which we call a \emph{partial homomorphism}. A vertex can be put at position $j$ to extend $Q$ to a homomorphism of $H$ exactly when it lies in a certain outneighbourhood or inneighbourhood of the image of each of the two neighbours of $j$. We call these two neighbourhoods the \emph{completion sets} $X_Q$ and $Y_Q$ of $Q$.
As each completion set has more than $N/2$ vertices, $X_Q\cap Y_Q\neq\emptyset$, so every partial homomorphism extends to a homomorphism of $H$. In particular, if $M$ is the maximum weight of $w$, then every partial homomorphism has cost at least $1-M$, and if a partial homomorphism has cost $c$, then every vertex of $X_Q\cap Y_Q$ has weight at least $1-c$.

The natural way to exploit this is to look at the vertices of maximum weight. Let $C$ be the set of vertices of weight $M$, and take an optimal fractional tiling. By complementary slackness, every homomorphism in this tiling has cost exactly one, and every vertex of $C$ is covered with total multiplicity one. Suppose, for simplicity, that each of these homomorphisms has at most one position mapped to $C$. Deleting this position from each homomorphism which meets $C$ gives a measure $\rho$ of total mass $|C|$ on partial homomorphisms into $F-C$, each of which costs $1-M$, the minimum possible cost.
Let $p(v)$ be the total weight, under $\rho$, of the positions mapped to $v$, so that $0\le p\le 1$ and $p$ has total $(r-1)|C|$.
Now take any partial homomorphism $Q$ of minimum cost $1-M$. Only vertices of weight $M$ can complete $Q$, so $X_Q\cap Y_Q\subseteq C$, and each vertex $v\notin C$ lies in at most one of $X_Q,Y_Q$, and so it contributes at most $1-p(v)+p(v)(\ind_{X_Q}(v)+\ind_{Y_Q}(v))$ to $|X_Q|+|Y_Q|$. Summing over all vertices gives
\[
 2\delta^0(F)\le|X_Q|+|Y_Q|\le N+|C|-(r-1)|C|+p(X_Q)+p(Y_Q).
\]
Therefore, if
\begin{equation}\label{ov:endpoint}
p(X_Q)+p(Y_Q)\leq (r-2)|C|,
\end{equation}
then $2\delta^0(F)\le N$, a contradiction (Lemma~\ref{general:max-deletion}).
Thus, much of our work is to find a partial homomorphism of minimum cost whose completion sets satisfy~\eqref{ov:endpoint}.
As $p$ has total $(r-1)|C|$, while $\rho$ has mass $|C|$, the condition~\eqref{ov:endpoint} says that the completion sets of $Q$ miss, on average, at least one position of each partial homomorphism in the support of $\rho$.

\subsection{Reduction to two families of orientations}\label{sec:sketch-reduction}
Recall that $k_H$ is the absolute value of the sum of the signs of $H$, that is, the difference between the numbers of forward and backward edges. As $H$ has even length, $k_H$ is even.
If $q$ divides the sum of the signs, then reducing the height walk of $H$ modulo $q$ gives a homomorphism from $H$ to the directed $q$-cycle. Indeed, each step goes from a height to the next height or from the next height back.
Composing such homomorphisms with a fractional tiling of $F$ by directed $q$-cycles, averaged over rotations, gives a fractional tiling by $H$ in which every vertex has the right multiplicity (Lemma~\ref{height400:quotient}).
More generally, we can cut the word of $H$ into pieces whose sign sums are divisible by four or by six, and glue together homomorphisms to the directed four-cycle and six-cycle (Lemma~\ref{height400:root-quotient}).
This handles every $H$ with $k_H=0$ or $k_H\geq 4$, and every $H$ with $k_H=2$ for which some proper cyclic interval of the word has sign sum at most $-4$ (Lemma~\ref{compact390:structural-reduction}).

For the remaining orientations, $k_H=2$, and, in each proper cyclic interval, the sign sum is at least $-3$. For these, no such cutting exists, as pieces with nonnegative sign sums divisible by four or six cannot have total sign sum two. The simplest example is the four-cycle $Q$, whose word $+++-$ has sign sum two.
By considering the last time the height walk visits each height (Lemma~\ref{height400:last-exits}), we show that, after rotating, the word of each remaining orientation has one of the following two forms (Lemma~\ref{compact390:structural-reduction}).
\begin{props}
\item $++W$, where $W$ is a nonempty closed height walk of span at most three, starting at any height in its range;\label{ov:family1}
\item $+U+V$, where $U$ and $V$ are closed height walks on $\{0,1,2,3\}$ which start and end at height zero and both visit height three.\label{ov:family2}
\end{props}
Thus, Theorem~\ref{compact390:finite} follows from fractional tilings by the directed four-cycle and six-cycle and by the orientations in~\ref{ov:family1} and~\ref{ov:family2}.

\subsection{Directed cycles and short orientations}\label{sec:sketch-small}
For directed four-cycles and six-cycles (Section~\ref{sec:directed}), and for the short orientations $Q$, $++-++-$, $D_{4,2}$ and $D_{5,3}$ (Section~\ref{sec:spantwo}), we follow Section~\ref{sec:sketch-covers} directly. For most of these, a short argument shows that the maximum weight $M$ exceeds $1/2$, so that a homomorphism of cost one cannot have two positions of weight $M$ (for the directed six-cycle, we need a separate argument when $1/3<M\le 1/2$). It then remains to find a partial homomorphism of minimum cost which satisfies~\eqref{ov:endpoint}.
Suppose there is none. We then modify partial homomorphisms of minimum cost by putting a vertex at the deleted position and deleting one of its two neighbouring positions instead. This keeps the cost minimum if the new vertex is cheap enough, and, as this new partial homomorphism must also violate~\eqref{ov:endpoint}, we learn something about the neighbourhoods of its new end vertices.
Repeating such modifications, we can move the deleted position around the cycle. In this process, the vertices of weight zero play a special role: after a preliminary rescaling, every relevant partial homomorphism meets them, and we track how the positions mapped to vertices of weight zero are arranged.
In Section~\ref{norm391:section}, we prove the general results of this kind that we need. For example, consider the vertices of weight zero with no inneighbour of weight zero, and those with no outneighbour of weight zero. We show that, if no suitable partial homomorphism exists, then both of these sets carry little multiplicity (Lemma~\ref{norm394:zero-extremes}).

For general orientations with $k_H=2$, the maximum weight need not exceed $1/2$, and a tight homomorphism may then have two positions of maximum weight. However, between these two positions there are about $r/2-2$ disjoint edges of $H$, each of which costs at least the minimum total weight $\kappa$ of an edge of $F$. Thus, the argument above applies whenever $(r/2-2)\kappa>1-2M$, and in this way we obtain the following bounds on an optimal cover $w$ with total weight less than $N/r$ (Theorem~\ref{norm391:normalized-cover}). The cover $w$ has a vertex of weight zero, and
\begin{equation}\label{ov:normalized}
 \frac 2r\le M\le\frac12,\qquad \left(\frac r2-2\right)\kappa\le 1-2M.
\end{equation}
Together with a simple counting argument (Lemma~\ref{norm391:savings}), this shows that some edge has total weight less than $1/r$, half the average total weight of an edge in a homomorphism of cost one.
For the orientations in~\ref{ov:family1} with span at most two, this is already enough (Section~\ref{sec:spantwo}). Indeed, for these we find a partial homomorphism which sends two consecutive heights, of large total multiplicity, to an edge of weight $\kappa$, and the other heights to cheap neighbours. Then \eqref{ov:normalized} shows that this partial homomorphism costs less than $1-M$, which is impossible (Lemma~\ref{height390:heavy-pair}).

\subsection{Walks of span three}\label{sec:sketch-spanthree}
For the orientations in~\ref{ov:family1} where $W$ has span three (Section~\ref{height390:three-section}) and those in~\ref{ov:family2} (Section~\ref{sec:tworeturns}), we use a different argument, which counts weights in neighbourhoods directly.
We fix an edge $x\to y$ of minimum total weight (or, in some cases, a vertex of weight zero), and consider the four neighbourhoods
\[
 X=N^-(x),\quad Y=N^+(y),\quad O=N^+(x),\quad I=N^-(y),
\]
each of which has more than $N/2$ vertices.
Now take a vertex $z$ lying in an intersection of some of these sets, say in $X\cap Y\cap O$. Then $x$, $y$ and $z$ span a directed triangle $x\to y\to z\to x$ together with the edge $x\to z$. We construct homomorphisms from $H$ into this small digraph by sending the heights of the height walk of $H$ to $x$, $y$ and $z$ in a suitable way. The cover condition for such a homomorphism then gives a lower bound on the weight of $z$, which depends on how many positions of $H$ are mapped to $z$.
These numbers of positions depend only on how often the height walks of $W$, or of $U$ and $V$, cross between consecutive heights. Thus, for each of the sixteen ways in which a vertex may lie in some of $X,Y,O,I$, we obtain an explicit lower bound on its weight in terms of these \emph{crossing counts} and of the weights of $x$ and $y$.

Finally, we combine these bounds as follows. Rescale the weights to $\omega=(r/2)w$, so that a cover of total weight less than $N/r$ has $\sum\omega<N/2$. Suppose that we find coefficients $c_X,c_Y,c_O,c_I\geq 0$ and $D\geq 0$ with $c_X+c_Y+c_O+c_I=1+2D$ such that
\[
 \omega(z)\geq c_X\ind_X(z)+c_Y\ind_Y(z)+c_O\ind_O(z)+c_I\ind_I(z)-D\qquad\text{for every vertex }z.
\]
Summing over all vertices and using that each of $X,Y,O,I$ has more than $N/2$ vertices gives $\sum\omega>(1+2D)N/2-DN=N/2$, a contradiction. Finding such coefficients amounts to checking sixteen inequalities, one for each pattern of memberships, against our lower bounds (see, for example, Lemma~\ref{root391:large-sum-count}).
Sometimes the lower bounds are too weak for vertices of small weight, and we supplement them with bounds on the number of such vertices, which come from the same kind of argument (Lemma~\ref{compact390:cuts}).
Most of Sections~\ref{height390:three-section} and~\ref{sec:tworeturns} consists of choosing the coefficients and checking these inequalities for all possible crossing counts.

Many of these checks are polynomial identities, or polynomial inequalities in the crossing counts on explicit integer ranges. We verify them by computer in exact rational arithmetic, and refer to each verification by a label such as~\chk{9.E}; the accompanying scripts (ancillary files \texttt{verify/*.py}) list the checks under the same labels. Each inequality is proved on its whole, infinite range. After a linear change of variables, the range becomes the nonnegative orthant, possibly with finitely many small cases removed and checked one by one. The polynomial then has nonnegative coefficients, either directly or after multiplication by a power of $1$ plus the sum of the variables (P\'olya's theorem). A condition that is affine in the weights is checked at the vertices of the relevant polygon or interval. No numerical approximation is used.

\subsection{Proof of Theorem~\ref{compact390:finite} subject to the results for each family}\label{sec:assembly}
The results described above are proved in Theorems~\ref{c4:thm:fractional}, \ref{c6:thm:fractional}, \ref{height390:three-width-factor} and~\ref{compact390:two-block}, and the reduction in Lemma~\ref{compact390:structural-reduction}. We now combine them.

\begin{proof}[Proof of Theorem~\ref{compact390:finite}]
By Theorems~\ref{c4:thm:fractional} and~\ref{c6:thm:fractional}, every finite digraph with minimum semidegree greater than half its order has perfect fractional homomorphic tilings by the directed four-cycle and by the directed six-cycle.
Therefore, by Lemma~\ref{compact390:structural-reduction}, it suffices to prove Theorem~\ref{compact390:finite} for the orientations in the two families~\ref{ov:family1} and~\ref{ov:family2}, and these are handled by
Theorems~\ref{height390:three-width-factor} and~\ref{compact390:two-block}, respectively.
\end{proof}
